% 编译方式：在本目录执行 xelatex main.tex 两次。
\documentclass[UTF8,a4paper,zihao=-4,fontset=windows]{ctexart}
\usepackage[left=20mm,right=20mm,top=22mm,bottom=22mm]{geometry}
\usepackage{amsmath,amssymb,booktabs,tabularx,longtable,array,graphicx,float,placeins}
\usepackage{caption,hyperref}
\hypersetup{hidelinks,pdftitle={考虑预测更新与电价不确定性的微网储能滚动调度}}
\captionsetup{font=small,labelfont=bf}
\setlength{\parindent}{2em}
\linespread{1.16}
\setlength{\emergencystretch}{2em}
\renewcommand{\arraystretch}{1.18}
\newcommand{\pos}[1]{\left[#1\right]_{+}}
\newcommand{\E}{\mathbb{E}}
\newcommand{\fig}[1]{\input{figures/#1.tex}}
\title{考虑预测更新与电价不确定性的微网储能滚动调度}
\author{}
\date{}
\begin{document}
\maketitle
\vspace{-2.8em}
\begin{abstract}
针对含光伏、负载与储能设备的微网购电调度问题，建立从典型日确定性优化到历史场景日前决策、预报更新日内调整及波动电价扩展的模型体系。统一采用十分钟离散时段，以储电量递推、容量和功率限制描述储能，并在明确的信息与结算条件下比较实际购电费用。

对于问题一，构建确定性线性规划，以混合整数模型及逐段约束检查作为对照。典型日购电费为35126.95元，较无储能基线节省12925.10元，降幅26.90\%。对于问题二，利用历史预测残差构造56天等权场景，在48小时前瞻窗口内优化日前购电和储能，逐日执行并结算紧急购电。2025年2—12月费用为14389135.08元；统一有效残差口径相对原方案减少12861.89元，改善集中于启动期。

对于问题三，将发布预报与历史预测在线融合，通过经验场景树表达未来信息更新，并在0、6、12、18点滚动决策。主方案费用为14022396.98元，较同口径历史预测对照节省367495.52元；增加18点更新未表现出明确的额外收益。对于问题四，日前联合价格场景方案和四次更新扩展方案的实际费用分别为15163043.89元与15061735.43元。通过同价结算对照、费用分解和逐段差值，分析价格信息及调度变化；同时披露4-3现有优化目标与实际结算之间的附加项差异。

结果表明，储能的经济作用取决于跨时段调节、预测误差与结算机制的共同作用，增加模型复杂度或更新次数并不保证本年度实际费用下降。进一步列示两种效率口径的已有费用对照，区分物理参数变化与主方案结果。本文结论限于已有历史回测，不将有限场景优化等同于全年真实随机控制的全局最优。

\noindent\textbf{关键词：}微网储能；随机规划；滚动优化；预报融合；电价场景
\end{abstract}
\clearpage
\section{问题重述与分析}
\subsection{研究对象与任务}
题目研究由光伏、负载、储能和外部电网构成的非孤岛式微网\cite{problem}。光伏供给与负载需求在时间上并不一致，电价又随时段变化。储能能够转移电量，但受到能量容量、充放电功率和转换损耗限制。购电不足时的紧急补给以及日内修改购电计划产生的费用，使“按预计缺口购电”不能直接等同于总费用最小的策略。

依据四问主模型及其正式结果\cite{models}，问题一讨论已知典型日负载、光伏与电价时的确定性调度；问题二在实际负载和光伏尚未实现时制定日前计划；问题三进一步利用0、6、12、18点发布的光伏预报调整决策；问题四分别在第二、三问框架中引入波动电价。本文保留五组既有主结果，不以其他历史候选替换主方案。

\subsection{建模思路与各问衔接}
第一问用于建立共同物理模型并解释储能收益。第二问的关键是平衡常规购电与高价紧急购电：过多计划购电会产生闲置费用，过少则增加缺口风险，因此以历史误差场景近似缺口的期望损失。跨日前瞻用于保留日末电量对次日的价值，避免人为逐日清空或复位电池。

第三问新增的信息并非实际未来光伏，而是逐次发布的预报。模型需要同时回答如何融合预报、如何表达未来信息分支，以及何时为修改计划支付费用。第四问的新增难点是电价与缺口可能相关，应在同一结算口径下分析费用变化。4-2继承日前决策，4-3继承多时点调整；二者的初始化与决策权限不同，不能仅凭总费用判断哪一算法更优。

\section{统一假设、数据口径与符号}
\subsection{共同约定与局部条件}
本文按以下约定建立主模型。
\begin{enumerate}
\item 每天划分为144个十分钟时段，段内功率视为平均功率。时间标签采用区间末时刻，00:10表示00:00—00:10，电量由功率乘$\Delta t=1/6$小时得到。
\item 不向外网售电；允许弃光。随机调度中允许弃置多余供给，已计划但未使用的购电仍计费。未加入折旧或寿命成本。
\item 题给额定容量12000 kWh、运行储电量1200—10800 kWh及功率上限5000 kW。主模型将功率限制解释为设备外端限制，并将“充放电效率90\%”解释为充电与放电效率各0.9，即往返0.81。另一效率口径在第\ref{sec:efficiency}节单独讨论。
\item 日前决策仅使用决策日前已实现的数据；日内决策可使用已发布预报和截至当时可见的信息。未来实际值仅用于事后结算，不用于选择当前控制。
\item 紧急购电按对应供电时段实际电价的5倍结算。第三问和4-3的取消量退回原费、另付50\%违约费，增购按1.5倍计价。该取消规则是本模型明确采用的结算解释。
\end{enumerate}

附件1对应典型日电价、负载和光伏预测曲线，第一问将该预测作为确定性输入；附件2对应全年实际负载和光伏，附件3对应发布预报，附件4对应实际波动电价。附件5规定五份结果文件及工作表内容。这里沿用主模型已有的时间转换、历史预测和样本处理，不额外假定完成了未记录的数据清洗。

\begin{table}[htbp]\centering\small
\caption{五组结果的统计期与储电量条件}\label{tab:conditions}
\begin{tabularx}{\linewidth}{l X r r}
\toprule 结果集 & 统计期及初始化 & 起始电量/kWh & 终止电量/kWh\\\midrule
问题一 & 典型日，日初固定 & 6000 & 6000\\
问题二、4-2 & 2—12月；1月共同调度预热 & 8550 & 6000\\
问题三、4-3 & 2—12月；1月保持电量并积累历史 & 6000 & 6000\\
\bottomrule\end{tabularx}\end{table}

附件5要求年度输出覆盖2025年2月1日至12月31日，共334天，1月用于主模型预热。题给2025年1月1日0点储电量6000 kWh；第一问另明确日初、日末相等。表\ref{tab:conditions}中的2月初状态及年度模型的年末6000 kWh约束是现有模型的初始化和终端约定，不冒充题面逐字要求。第三问使用统一6000 kWh起点的历史预测对照，而非直接把第二问正式费用作为其同条件基准。一般滚动窗口末端条件也不同于每天日末回到日初。

题面还要求在论文中按表1、表2、表3给出指定结果，并讨论是否需要其他时刻的预报。本文在各问呈现总费用与图示，在附录\ref{sec:outputs}完整列示指定十分钟购电、四小时充放电和连续紧急购电区间；日期为3月20日、6月21日、9月23日及12月21日。第三问的更新频率比较直接回应预报时刻选择。第四问按题意重算第二、三问，不额外设立4-1。

\subsection{符号与储能约束}
\begin{table}[htbp]\centering\small
\caption{主要符号与单位}
\begin{tabularx}{\linewidth}{l X l}\toprule
符号 & 含义 & 单位\\\midrule
$d,t,s,n$ & 日期、日内时段、历史场景、场景树节点索引 & —\\
$L_t,V_t,N_t$ & 负载、光伏及净负载电量，$N_t=L_t-V_t$ & kWh\\
$\widehat N_t,N_{s,t}$ & 净负载点预测与场景值 & kWh\\
$g_t,g_t^0,q_t$ & 日前购电、0点原计划、最终调整后购电 & kWh\\
$c_t,b_t,E_t$ & 外端充电量、外端放电量、段末内部储电量 & kWh\\
$e_t,w_t$ & 紧急购电量、弃光或弃置供给量 & kWh\\
$u_t,v_t$ & 相对原计划的增购量与取消量 & kWh\\
$p_t,\widehat p_t,p_{s,t}$ & 实际、预测与场景电价 & 元/kWh\\
$\omega_s,\pi_n$ & 场景权重与节点概率 & —\\
$\eta_c,\eta_d$ & 充电效率、放电效率 & —\\
\bottomrule\end{tabularx}\end{table}

为避免日期索引$d$与放电量混淆，全文用$b_t$表示主程序中的放电变量。共同储能约束写为
\begin{align}
E_t&=E_{t-1}+\eta_c c_t-b_t/\eta_d,\label{eq:storage}\\
1200\le E_t&\le10800,\qquad 0\le c_t,b_t\le5000/6.\label{eq:bounds}
\end{align}
式\eqref{eq:storage}区分设备外端电量与内部库存，不能直接把充电量减放电量作为库存变化。跨日时下一天的起点必须等于上一天实际执行后的末电量。

\section{问题一：典型日确定性储能调度}
\subsection{模型建立与求解}
给定全天负载、光伏和电价，以外网购电费为目标：
\begin{align}
\min\quad &J_1=\sum_{t=1}^{144}p_tg_t,\label{eq:q1obj}\\
\text{s.t.}\quad &g_t+V_t+b_t=L_t+c_t+w_t,\label{eq:q1balance}\\
&g_t\ge0,\qquad 0\le w_t\le V_t,\\
&E_0=E_{144}=6000,
\end{align}
并满足式\eqref{eq:storage}—\eqref{eq:bounds}。弃光变量使光伏富余时仍可保持能量平衡，首末状态相同则排除了消耗初始库存获得表面节省的情况。

模型首先按线性规划求解，再用二元变量$y_t$构造互斥对照：$c_t\le(5000/6)y_t$、$b_t\le(5000/6)(1-y_t)$。已有结果显示，本实例LP解的144个时段均不存在同时充放电，且LP与MILP目标相同。因此，对本实例而言，LP解已经是互斥模型的可行最优解。这里采用计算结果与约束核验作为依据，不把缺少边界条件的扰动论证推广到所有含弃光上界的模型。

\subsection{最优调度与费用}
\begin{table}[htbp]\centering\small
\caption{典型日储能调度与无储能基线}
\begin{tabular}{lrr}\toprule
方案 & 购电费/元 & 较基线节省/元\\\midrule
无储能 & 48052.05 & —\\
储能主方案 & 35126.95 & 12925.10\\
\bottomrule\end{tabular}\end{table}

无储能基线逐段优先利用光伏，不足部分由外网补齐，即$g_t^{\mathrm{base}}=\max(L_t-V_t,0)$。主方案购电费35126.95元，相对48052.05元基线下降26.90\%。调度过程中既有低价购电后储存，也有吸收光伏富余后在高价时段放电，经济作用来自全日约束下的联合安排。

\fig{01_q1_dispatch}
图\ref{fig:01_q1_dispatch}按同一时间轴对齐电价、供需、充放电与库存。储电量在允许区间内变化，日初和日末均为6000 kWh。低价时段充电与高价时段放电是总体特征，但局部行为还受光伏、负载、库存和功率限制影响，不能仅按电价阈值解释每个时段。

\subsection{费用变化的恒等分解}
为避免混淆效益归因，按逐段购电量变化定义
\begin{align}
A&=\sum_t p_t\max(g_t^{\mathrm{base}}-g_t,0),\\
B&=\sum_t p_t\max(g_t-g_t^{\mathrm{base}},0),\\
J_1&=J_1^{\mathrm{base}}-A+B.
\end{align}
少购电减少费用19111.64元，其他时段增购增加费用6186.54元，抵消后节省12925.10元。该分解直接由输入和调度复算，是现金费用的恒等关系，不声称把光伏消纳与套利的因果贡献完全分离。
\fig{02_q1_savings}
图\ref{fig:02_q1_savings}说明，部分时段增购电与全天节省并不矛盾：储能使电量从相对低成本时段流向高成本时段。完整逐段结果保留在正式交付表，正文给出题目指定的六个购电时段及四小时充放电汇总，见附录\ref{sec:outputs}。
\FloatBarrier

\section{问题二：历史误差场景下的滚动日前调度}
\subsection{预测与有效残差样本}
日前决策不能使用当天实际负载和光伏。负载预测采用最近4次同星期几的均值，并叠加最近5天预测残差的均值；光伏预测采用最近7天均值和最近28天残差漂移。将两者转换为电量后相减，得到净负载点预测$\widehat N_{d,t}$。所有残差均以该历史日当时可得的预测计算，而不是用事后拟合值替代。

场景集合取最近56天内的有效历史日$\mathcal H_d$。对$r_s\in\mathcal H_d$，
\begin{equation}
N_{s,t}=\widehat N_{d,t}+N_{r_s,t}-\widehat N_{r_s,t},\qquad
\omega_s=1/|\mathcal H_d|.\label{eq:residual}
\end{equation}
前7天缺少同星期历史，负载预测使用典型日回退，因此最终方案从1月8日起纳入残差。早期回退误差本身是合法的因果预测误差，排除它们是为了使场景样本与后续预测器口径更一致。2月1日可用24条残差，随后逐渐增至56条。主模型原有预测处理保持不变，不把其他预测器改动混入该对照。

\subsection{两阶段缺口补救与48小时前瞻}
普通日采用288个时段的窗口。当天购电和充放电是场景间共享的决策，紧急购电$e_{s,t}$作为场景相关的补救变量；次日使用点预测近似。优化目标为
\begin{equation}
\min\ J_2^{\mathrm{look}}=
\sum_{k=1}^{288}p_k g_k+
5\sum_s\omega_s\sum_{t=1}^{144}p_t e_{s,t}.\label{eq:q2}
\end{equation}
其中$p_k$为重复至两日的典型日电价。需求约束为
\begin{align}
g_k+b_k-c_k&\ge\widehat N_k,\qquad g_k\ge0,\\
g_t+b_t-c_t+e_{s,t}&\ge N_{s,t},\qquad e_{s,t}\ge0.
\end{align}
储能满足共同约束，普通日窗口末端回到本次窗口起始储电量。次日前瞻用于估计库存价值，不代表次日计划已成交；当天执行完毕后重新计算后续计划。12月31日仅保留当天，并固定年末电量6000 kWh。

在固定充放电、单时段正电价且没有其他约束时，$pg+5p\E[(N-g-b+c)_+]$的内部最优条件对应80\%分位数。这有助于解释高价缺口补救为何推动计划购电增加，但完整模型还含储能耦合、点预测供给约束和终端条件，不能逐段直接套用该分位数替代求解。

\subsection{执行流程与物理可行性}
每天首先更新预测与历史场景，求解稀疏LP，然后只执行当天144段，传递实际日末储电量。实际缺口与费用分别为
\begin{equation}
e_t=\max(N_t-g_t-b_t+c_t,0),\qquad
J_{2,d}=\sum_t p_tg_t+5\sum_t p_te_t.\label{eq:q2settle}
\end{equation}
计划购电费与未用电量无关，故降低紧急购电量并非唯一优化目标。

对于允许弃置多余供给的随机模型，设$\rho=\eta_c\eta_d<1$。如出现同时充放电，取$a=\min(c_t,b_t/\rho)$并令$c'_t=c_t-a$、$b'_t=b_t-\rho a$。该变换保持式\eqref{eq:storage}不变，净供给增加$(1-\rho)a$，功率不增，且至少一个充放电量变为零。允许弃置时可保持可行，紧急购电不会增加。主程序采用这一等库存处理并核验执行结果。该推理依赖相应弃置条件，不能无条件推广到其他设备模型。

\subsection{结果与窗口规则比较}
\begin{table}[htbp]\centering\small
\caption{第二问两种残差口径下的费用（元）}
\begin{tabular}{lrr}\toprule
费用项 & 含回退残差 & 有效残差主方案\\\midrule
计划费 & 13508190.55 & 13462061.30\\
紧急费 & 893806.43 & 927073.78\\
总费用 & 14401996.98 & 14389135.08\\
\bottomrule\end{tabular}\end{table}

主方案总费用为14389135.08元。相对含回退残差的56天等权方案，计划费减少46129.24元，紧急费增加33267.35元，总费用净降12861.89元。表中各金额分别由未舍入值计算后保留两位小数，个别差额与已舍入单元格相减可差0.01元。

\fig{03_q2_residuals}
图\ref{fig:03_q2_residuals}采用相同规则在两种残差口径下的配对结果。在有效残差条件下，56天等权费用低于半衰期28天与四候选在线选择；据此保留较简单的等权方案。该比较是同年度回顾性分析，不能证明56天在其他年份仍最优。

从时间分布看，2月节省15521.02元，3月增加2659.13元，最后一个日费用有差异的日期为3月4日。3月5日至12月31日逐日费用一致，说明改动主要作用于启动阶段的场景构成，而非全年持续提高适应性。指定日期结果见附录\ref{sec:outputs}。
\FloatBarrier

\section{问题三：预报融合与多阶段滚动调整}
\subsection{发布预报的处理和在线融合}
附件3提供每次发布后24小时的整点光伏预报。先线性插值至十分钟网格，发布瞬间端点采用当时已观测光伏，0点采用上一日末观测。超过预报覆盖范围的次日时段用历史预测补足。负载仍由历史模型预测，不以当天尚未发生的实际负载更新。

设同一发布时刻、同一未来六小时块的历史预测为$H$、发布预报为$F$、实际光伏为$V$。采用最近56天内已经完整验证的预报样本，计算
\begin{equation}
\alpha=\operatorname{clip}_{[0,1]}
\frac{\sum(F-H)(V-H)}{\sum(F-H)^2}.\label{eq:fusion}
\end{equation}
以$\alpha F+(1-\alpha)H$为基础，再加历史平均残差乘$n/(n+14)$的收缩修正，并裁剪至非负。分母退化时按原实现回退；有效样本少于7条时暂用发布预报。每个发布时刻与未来六小时块分别估计，避免把夜间和日间的误差结构混为一谈。

跨午夜预报行只有在未来24小时全部实现后才能训练。因此，决策日使用的融合训练行最晚发布于日前两天；当天后续预报不能反向参与当前权重估计。56天窗口和收缩常数14沿用既有参数，不声称经过独立测试选型。

\fig{05_q3_forecast_fusion}
图\ref{fig:05_q3_forecast_fusion}显示，6—12时段的MAE由历史预测309.24 kW降至融合预测169.41 kW，12—18时段由288.99 kW降至133.68 kW。日夜面板使用不同纵轴尺度。融合包含偏差修正，不能仅将其理解为两条曲线的固定凸组合；预测误差改善也不直接等于同比例费用下降。

\subsection{场景树与非前视决策}
采用与式\eqref{eq:residual}同类的整日残差场景，但点预测和历史残差都对应当前发布时刻。每个历史日还携带后续预报修订向量，用确定性二叉聚类近似下一次发布可能带来的信息，最多形成四层、15个决策节点。聚类依据是届时可见的预报修订，不是尚未观测的实际负载或实际光伏。

设节点$n$包含历史场景集合$\mathcal S_n$，概率为$\pi_n=|\mathcal S_n|/S$。同一节点只设置一组购电、充放电和库存变量，各场景仅缺口补救不同。因此，无法用当前信息区分的场景自动共享决策，满足非前视要求。子节点的初始储电量等于父节点末值；仍让全部历史样本参与缺口期望，不以少量均值曲线替代全部误差样本。

\subsection{调整合同与优化模型}
设$g_t^0$为0点原计划，$q_t$为最终执行购电，
\begin{equation}
u_t=\max(q_t-g_t^0,0),\qquad v_t=\max(g_t^0-q_t,0).
\end{equation}
采用取消退原费并另付50\%违约费的解释，则实际费用为
\begin{align}
J_3&=\sum_t\left(p_tg_t^0+1.5p_tu_t-p_tv_t+0.5p_tv_t+5p_te_t\right)\notag\\
&=\sum_t\left(p_tg_t^0+1.5p_tu_t-0.5p_tv_t+5p_te_t\right).\label{eq:q3settle}
\end{align}
退款不是负购电，净调整费可以为负。原计划费、净调整费和紧急费相加时，不再重复计入退款、违约费与增购费。题面使用“交易时刻电价”，现有模型按购电所属供电时段定价；如改为决策发布时刻定价，需改变优化与结算，不能只调整文字。

0点原计划跨场景共享，节点变量满足$q_{n,t}=g_t^0+u_{n,t}-v_{n,t}$以及
\begin{equation}
q_{n,t}+b_{n,t}-c_{n,t}+e_{s,t}\ge N_{s,t},\qquad s\in\mathcal S_n.
\end{equation}
另保留净计划供给不低于节点点预测的约束和共同储能条件。省略极小的吞吐量消歧项，0点优化目标为
\begin{equation}
\min\ \sum_t p_tg_t^0+
\sum_n\pi_n\sum_{t\in\mathcal T_n}(1.5p_tu_{n,t}-0.5p_tv_{n,t})
+\frac5S\sum_s\sum_t p_te_{s,t}+\Phi.\label{eq:q3obj}
\end{equation}
其中$\mathcal T_n$为节点执行区间，$\Phi$为各叶节点接入的次日确定性前瞻费用。0点第一执行块不调整，后续更新时原计划固定，原计划费作为常数从优化中省略，但实际结算仍完整保留。

普通日的次日末端固定为本日0点储电量，6、12、18点重算时不随当前库存改变该终端目标。12月31日直接固定当日末6000 kWh。每次仅执行本发布时刻至下一允许更新时刻的控制，其后计划只是预案；每个供电时段最多实际调整一次，不依赖反复买卖后仅结算净差额的假设。

\subsection{更新频率的经济价值}
\begin{table}[htbp]\centering\small
\caption{第三问更新频率对照}
\begin{tabular}{lrr}\toprule
方案 & 总费用/元 & 紧急购电量/kWh\\\midrule
同口径历史预测 & 14389892.50 & 246099.33\\
仅0点融合 & 14181445.03 & 218589.30\\
0、6点 & 14101371.47 & 186570.89\\
0、6、12点 & 14022358.10 & 175344.39\\
四次更新主方案 & 14022396.98 & 175369.95\\
\bottomrule\end{tabular}\end{table}

同口径历史预测对照为14389892.50元，其初始化和物理处理与本问一致，并非第二问正式结果14389135.08元。四次更新主方案总费用14022396.98元，节省367495.52元，紧急购电量由246099.33 kWh降至175369.95 kWh。

\fig{04_q3_information_value}
图\ref{fig:04_q3_information_value}把改善分为引入0点融合和后续新增调整。仅0点融合节省208447.47元；增加6点、12点分别再节省80073.56元和79013.36元；加入18点后费用反增38.88元。故本年度数据支持重视6点和12点调整，未显示18点的明确额外经济收益，预设四次更新主结果仍予保留。

这不意味着18点信息始终无价值。它可能改变次日预报与库存安排，且本方法在有限历史样本上重建树并滚动执行，真实费用并不保证随更新机会增加而下降。不能利用同年最好的更新组合，反向宣称它是事先已知的最优规则。

\begin{table}[htbp]\centering\small
\caption{第三问主结果实际费用分解}
\begin{tabular}{lr}\toprule
项目 & 费用/元\\\midrule
原计划费 & 13345819.30\\
增购费 & 89947.28\\
取消退款（扣减） & -150170.13\\
违约费 & 75085.07\\
净调整费 & 14862.22\\
紧急费 & 661715.46\\
总费用 & 14022396.98\\
\bottomrule\end{tabular}\end{table}

\subsection{指定日期的日内行为}
\fig{06_q3_rolling_adjustment}
图\ref{fig:06_q3_rolling_adjustment}展示四个指定日期最终执行购电减0点原计划的功率差，并对齐储电量和紧急购电时段。正值表示增购，负值表示取消；橙色短线只标记紧急购电发生时间，不表示数量。该图不是全部中间计划的修订历史。库存轨迹说明，日内调整既影响当前缺口，也受到后续用电和终端条件制约。
\FloatBarrier

\section{问题四：波动电价下的购电调度}
\subsection{4-2：价格与净负载联合场景}
题面给出了实际价格，但未明确每天0点是否已知全天价格。因此主模型采用“当天价格未知、仅有历史观测”的信息条件；另设当日电价已知的对照。实际价格用于事后结算，不能因附件包含全年价格而提前使用未来值。

以1月8—31日的24个历史预测日比较前一日、近7日均值、近7日形状加前一日水平、最近4次同星期几均值四种规则，按预测RMSE选定最后一种，并从2月起固定。样本日期必须早于决策日，预测次日时也不能使用尚未实现的当日价格。缺少同星期样本时用近7日均值回退。

将与净负载残差同日的电价预测误差配对，得到
\begin{equation}
p_{s,t}=\max\{0.001,\widehat p_{d,t}+p_{r_s,t}-\widehat p_{r_s,t}\},\qquad
\overline p_{d,t}=\sum_s\omega_s p_{s,t}.
\end{equation}
0.001元/kWh为场景正值下限，不改变原始实际价格。模型保持第二问的共享日前控制、56天等权及48小时前瞻，目标改为
\begin{equation}
\min\ \sum_t\overline p_{d,t}g_t+
5\sum_s\omega_s\sum_t p_{s,t}e_{s,t}+
\sum_t\widehat p_{d+1|d,t}g_{144+t}.\label{eq:q42}
\end{equation}
对于净供给$z_t=g_t+b_t-c_t$，一般有
$\E[p_t(N_t-z_t)_+]\ne\E[p_t]\E[(N_t-z_t)_+]$，故价格与缺口的依赖关系可能影响风险权重。这里直接相关的是同一时段的联合分布，保留整日配对本身不能保证额外收益。均价对照使用相同净负载场景，将所有价格场景替换为其均值，以检验该关系的实际作用。

\begin{table}[htbp]\centering\small
\caption{4-2同价结算下的价格信息对照}
\begin{tabular}{lrr}\toprule
方案 & 实际费用/元 & 较重结算基线节省/元\\\midrule
原Q2计划重结算 & 15184704.23 & 0.00\\
均价对照 & 15160921.45 & 23782.79\\
联合场景主方案 & 15163043.89 & 21660.34\\
当日电价已知对照 & 15033043.18 & 151661.05\\
\bottomrule\end{tabular}\end{table}

\fig{07_q42_price_information}
联合场景主方案费用15163043.89元，相对第二问原计划按附件4实际价格重结算的15184704.23元节省21660.34元，约0.143\%。均价对照费用15160921.45元，反而比联合方案低2122.44元。因此，本年结果支持存在计划费与紧急费的权衡，不能据此声称联合场景必然更优。当日电价已知对照为15033043.18元，放宽了信息条件，但仍受既有预测和滚动近似限制，不是已证明的理论最优下界。

\fig{08_q42_schedule_shift}
图\ref{fig:08_q42_schedule_shift}对齐实际电价、场景均价与主方案相对重结算基线的日前购电变化。两条价格曲线间的填色是差值提示，不是置信区间。购电变化还受到净负载场景、储能状态和次日前瞻影响，不能把每个变化都归因于相邻时段的价格偏差。
\FloatBarrier

\subsection{4-3：随机价格下的日内调整与现有实现边界}
沿用第三问的融合、场景树及调整合同，未来电价预测采用此前7天同一时段均值，已执行时段使用已知价格。电价误差与净负载误差来自相同历史日，场景正值下限为$10^{-5}$元/kWh。其价格预测规则与4-2不同，应分别说明，不能合写成同一个电价预测器。

\subsubsection{优化目标与实际结算的区别}\label{sec:q43objective}
现有主程序除0点原计划价格项、增购/取消项、紧急补救项及次日前瞻外，还给节点调整后购电$q_{n,t}$设置了价格系数。令$\widetilde p_t$为当前历史价格预测，$\overline p_{n,t}$为节点内价格场景均值，则0点实际求解的目标可表示为
\begin{align}
\widetilde J_{43}={}&\sum_t\widetilde p_t g_t^0
+\sum_n\pi_n\sum_{t\in\mathcal T_n}\overline p_{n,t}q_{n,t}\notag\\
&+\sum_n\pi_n\sum_{t\in\mathcal T_n}\overline p_{n,t}(1.5u_{n,t}-0.5v_{n,t})\notag\\
&+\frac5S\sum_s\sum_t p_{s,t}e_{s,t}+\widetilde\Phi.\label{eq:q43actual}
\end{align}
后续更新时原计划固定，其首项省略为常数。现有次日前瞻使用当前价格预测向量，并按对应节点概率加权；它是原实现的近似，非次日价格的完美信息。

式\eqref{eq:q43actual}中的$\sum\pi_n\overline p_{n,t}q_{n,t}$在第三问优化中不存在，也不属于既定合同实际费用的额外应付款。实际报告仍按式\eqref{eq:q3settle}代入附件4价格计算。因而当前4-3结果是这一既有决策目标生成、按合同重新结算的策略结果，不能表述为对合同期望总费用目标的完全一致求解。该附加项可能影响购电和库存，尚不能在未重新优化对照的情况下量化其影响；本文保留正式结果并将此列为实质性待核对问题，不为附加项补造未经记录的经济动机。

\subsubsection{实际费用及与第三问的差异}
\begin{table}[htbp]\centering\small
\caption{4-3主结果实际费用分解}
\begin{tabular}{lr}\toprule
项目 & 费用/元\\\midrule
原计划费 & 13457230.68\\
增购费 & 152250.38\\
取消退款（扣减） & -173200.59\\
违约费 & 86600.30\\
净调整费 & 65650.08\\
紧急费 & 1538854.67\\
总费用 & 15061735.43\\
\bottomrule\end{tabular}\end{table}

4-3实际总费用15061735.43元，由原计划费13457230.68元、净调整费65650.08元和紧急费1538854.67元组成。相对第三问增加1039338.45元，其中紧急费增加877139.21元，占账面差额84.39\%。该比例只说明费用分项，不证明额外成本全部由电价波动或预报误差造成。

\fig{09_q43_incremental_comparison}
图\ref{fig:09_q43_incremental_comparison}同时区分跨情形费用差额与情形内部新增更新的效果。4-3从本问历史预测对照到四次更新，共减少455249.11元；其中18点相对0、6、12点方案减少6875.04元。它与第三问18点增加38.88元的结果不同，进一步说明更新时点价值依赖价格、预测与决策目标，不能把固定电价下的频率结论直接移植过来。

\begin{table}[htbp]\centering\small
\caption{4-3更新频率对照}
\begin{tabular}{lrr}\toprule
方案 & 总费用/元 & 紧急购电量/kWh\\\midrule
同口径历史预测 & 15516984.54 & 504526.13\\
仅0点融合 & 15295120.10 & 461933.56\\
0、6点 & 15173094.46 & 410511.06\\
0、6、12点 & 15068610.47 & 387077.17\\
四次更新主方案 & 15061735.43 & 383981.09\\
\bottomrule\end{tabular}\end{table}

\fig{10_q43_schedule_differences}
图\ref{fig:10_q43_schedule_differences}按同日同段将4-3减第三问，分别展示价格、最终购电功率和段末储电量差值。每个面板保留四日共576个十分钟值，使用以零为中心的全范围对称色标。两问年度初末状态相同，但指定日初库存可能已不同，因此库存差含前期历史路径影响。此处还有式\eqref{eq:q43actual}所示目标差异，图中共同变化不能用作单一价格因素的因果识别。
\FloatBarrier

\section{模型检验与效率口径对照}
\subsection{主结果的物理和费用核验}
主模型已有验证记录包含储电量递推、容量与功率上下界、充放电互斥、跨日连续性、紧急缺口、逐段结算、日汇总及结果表反读。第一问共144段，每个年度结果集分别包含334天、48096段，不能将48096写成全部四问合计。

第一问记录的能量平衡最大残差约$2.27\times10^{-13}$ kWh，储电量递推最大残差约$7.28\times10^{-12}$ kWh。第三问与4-3均无执行充放电重叠，储电量递推最大误差约$4.55\times10^{-12}$ kWh，实际费用分别复算为14022396.98元和15061735.43元。这里引用既有记录，不将此次论文整理描述为重新运行全部求解器。

费用复算能确认给定调度的账目正确，却不能单独证明优化目标正确。第\ref{sec:q43objective}节所示4-3目标附加项不会因实际费用核验通过而消失；该问题需将原目标与修正目标进行独立求解对照后才能量化影响。

\subsection{信息边界与结算解释}
已有程序通过扰动未来实际值及尚未发布预报检查当前预测与决策是否改变；这类检查针对执行阶段的前视读取。它不能证明原预测参数在历史研发中从未参考测试年，因此本文将结论限定为给定参数下的历史滚动回测，不称为完全独立的外部验证。

取消购电的退款规则会影响实际费用。按原主调度不变、取消不退原费且仍另付50\%违约费的解释，第三问费用为14172567.11元，4-3为15234936.02元。这是同一调度的结算敏感性，不是另一合同下重新优化的最优费用。正文主结果始终采用式\eqref{eq:q3settle}，不得将两种结算混用。

\subsection{储能效率口径对结果的影响}\label{sec:efficiency}
题面仅给出“充放电效率90\%”，没有进一步区分单向与往返。主模型按单向各0.9建立；已有补充材料\cite{efficiency}另采用往返0.9并对称拆分损耗。两者的物理递推分别为
\begin{align}
\text{主口径：}\quad&E_t=E_{t-1}+0.9c_t-b_t/0.9,\\
\text{对照口径：}\quad&E_t=E_{t-1}+\eta c_t-b_t/\eta,\quad\eta=\sqrt{0.9}.
\end{align}
对称拆分是对照模型的明确选择，不意味着相同往返效率下所有单向拆分都会产生完全相同的受限调度。

\begin{table}[htbp]\centering\small
\caption{已有两种效率口径总费用对照（元）}\label{tab:efficiency}
\begin{tabular}{lrrrr}\toprule
结果集 & 往返0.81 & 往返0.90 & 差额 & 变化率\\\midrule
问题一 & 35126.95 & 33801.50 & -1325.45 & -3.77\%\\
问题二 & 14389135.08 & 13910066.83 & -479068.25 & -3.33\%\\
问题三 & 14022396.98 & 13558297.17 & -464099.81 & -3.31\%\\
问题四-2 & 15163043.89 & 14684028.55 & -479015.34 & -3.16\%\\
问题四-3 & 15061735.43 & 14579714.65 & -482020.78 & -3.20\%\\
\bottomrule\end{tabular}\end{table}

表\ref{tab:efficiency}逐项列示补充材料已经报告的总费用。对照口径下五组费用下降约3.16\%—3.77\%，说明效率解释对费用绝对值有可观影响，应在报告结果时声明，不能视为可忽略的数值误差。降低损耗会改变可用输出电量与库存安排，因此这种变化在机理上合理。

该表只有五组主结果的总费用，尚不足以证明所有更新频率排序、均价与联合场景排序在两种效率下都保持不变。本文也不以总费用差额断言某一独立套利或光伏消纳效应。效率对照采用补充材料的既有记录，本次未读取各问效率分支或重新核验其逐段输出，故不新增“两口径全部逐段核验通过”的结论。若用于最终定稿，应将相关论断限制在现有表格可支持的范围。

\subsection{图表与解释一致性}
所有主结果图取自新图包\cite{figures}，保持统一字体、配色与单位，并保留来源对应关系。差值图明确相减方向，费用图区分元和万元，紧急时段标记不冒充电量；热力图不截断极端值。第四问跨情形费用比较同时伴随模型条件变化，属于描述性比较。图区分了“变化是什么”与“变化由何导致”，避免用视觉相似性或相关性替代因果证据。

\section{模型评价}
\subsection{优点与适用范围}
共同的储能约束使五组结果能够以一致的物理单位解释；历史场景和滚动执行将预测风险、购电费用与库存价值联系起来。第三问通过节点共享变量表达信息限制，能够在不预知未来实际值的条件下利用预报更新。费用分项与同价结算对照使节省来源和比较口径可追溯，指定日期图表又提供了对全年汇总的具体解释。

\subsection{局限与改进方向}
首先，4-3当前目标含合同费用之外的附加项，现有结果不能证明针对合同总费用的最优性。优先改进应是核对目标系数与合同解释，再决定是否重算，不应直接将该项事后包装成风险偏好。

其次，经验场景来自有限历史窗口，小型二叉树和次日确定性近似无法覆盖全部随机结构。窗口、预测器及融合收缩参数虽在执行时固定，其跨年适用性仍需独立数据验证。第二问启动期改善和第三问18点微小差异，均不宜被放大为普遍规律。

最后，效率、功率作用位置、取消退款和定价时刻均存在明确建模解释；储能寿命、设备退化和售电机制不在现有主模型内。后续若增加这些因素，需要同步修改状态或成本，而非沿用现有费用直接外推。

\section{结论}
典型日确定性模型在首末库存相同的条件下得到35126.95元购电费，较无储能节省26.90\%，展示了储能协调时段供需与电价的经济作用。历史场景日前模型的年度费用为14389135.08元，有效残差筛选带来的12861.89元改善主要集中于启动期。

预报融合与场景树滚动主方案的费用为14022396.98元，较其同口径历史预测对照减少367495.52元。6点与12点更新在本年度具有较明确的增量收益，18点在固定电价主模型中未表现出额外节省。波动电价下，4-2主模型实际费用15163043.89元，均价对照略低，说明价格与缺口关系的建模价值必须由具体回测评估。

4-3既有策略的实际费用为15061735.43元，与第三问的费用及调度差异已经分项呈现，但其优化目标存在额外购电价格项，故当前结果应作为具有该实现条件的策略输出保留。两种效率口径的已有对照进一步表明，物理损耗假设会改变费用数值。以上结论共同强调：比较购电策略前，应明确物理、信息与结算条件，并把可核算的结果与尚待验证的解释分开。

\begin{thebibliography}{9}
\bibitem{problem} 2026年高教社杯全国大学生数学建模竞赛. C题：微网与外部电网电力调控策略[题目及附录]. data/C题.pdf, 2026.
\bibitem{models} 项目组. 四问主模型建模说明、程序及正式结果[内部研究材料]. question1—question4, 2026.
\bibitem{figures} 项目组. 新论文图包：图注与正文、绘图数据及来源清单[内部研究材料]. paper/figure\_package, 2026.
\bibitem{efficiency} 项目组. 储能效率口径的稳健性检验及说明[内部研究材料]. paper/extra, 2026.
\end{thebibliography}
\appendix
\section{指定时段与指定日期的完整结果}\label{sec:outputs}
以下表格按照题目表1、表2、表3所要求的字段组织。购电量、充放电量、储电量均以kWh计，费用以元计。第一问采用三组购电时段并列、两组充放电时段并列的原表布局；年度结果对每个指定日期分别列示。第三问和4-3同时给出0点原计划与最终调整后购电，充放电为实际执行量。紧急购电按连续发生区间汇总，四个日期并列，不只列每日合计。

\subsection{问题一指定结果}
\begin{table}[htbp]\centering\scriptsize
\caption{问题一：指定时段及全天购电}
\begin{tabular}{lrrrrr}\toprule
时间段 & 购电量 & 时间段 & 购电量 & 时间段 & 购电量\\\midrule
10:00--10:10 & 0.00 & 12:00--12:10 & 480.41 & 14:00--14:10 & 0.00\\
16:00--16:10 & 445.43 & 18:00--18:10 & 531.89 & 20:00--20:10 & 0.00\\
\multicolumn{2}{l}{全天计划购电量} & 59482.70 & \multicolumn{2}{l}{全天实际购电费} & 35126.95\\
\bottomrule\end{tabular}\end{table}

\begin{table}[htbp]\centering\scriptsize
\caption{问题一：四小时充放电与首末储电量}
\begin{tabular}{lrrrrr}\toprule
时间段 & 充电量 & 放电量 & 时间段 & 充电量 & 放电量\\\midrule
00:00--04:00 & 4500.00 & 0.00 & 04:00--08:00 & 833.33 & 6365.84\\
08:00--12:00 & 4787.96 & 1703.00 & 12:00--16:00 & 5286.04 & 91.10\\
16:00--20:00 & 0.00 & 5780.13 & 20:00--24:00 & 5333.33 & 2859.87\\
\multicolumn{2}{l}{0:00储电量} & 6000.00 & \multicolumn{2}{l}{24:00储电量} & 6000.00\\
\bottomrule\end{tabular}\end{table}

\FloatBarrier
\subsection{问题二指定日期结果}
\begin{table}[htbp]\centering\scriptsize
\caption{问题二 2025-03-20：指定时段及全天购电}
\begin{tabular}{lrrrrr}\toprule
时间段 & 购电量 & 时间段 & 购电量 & 时间段 & 购电量\\\midrule
10:00--10:10 & 0.00 & 12:00--12:10 & 626.63 & 14:00--14:10 & 0.00\\
16:00--16:10 & 494.24 & 18:00--18:10 & 706.68 & 20:00--20:10 & 0.00\\
\multicolumn{2}{l}{全天计划购电量} & 65615.94 & \multicolumn{2}{l}{全天实际购电费} & 43269.00\\
\bottomrule\end{tabular}\end{table}

\begin{table}[htbp]\centering\scriptsize
\caption{问题二 2025-03-20：四小时充放电与首末储电量}
\begin{tabular}{lrrrrr}\toprule
时间段 & 充电量 & 放电量 & 时间段 & 充电量 & 放电量\\\midrule
00:00--04:00 & 1666.67 & 0.00 & 04:00--08:00 & 833.33 & 5501.23\\
08:00--12:00 & 5912.26 & 3138.77 & 12:00--16:00 & 5454.26 & 566.88\\
16:00--20:00 & 0.00 & 5671.59 & 20:00--24:00 & 3498.06 & 2968.41\\
\multicolumn{2}{l}{0:00储电量} & 8550.00 & \multicolumn{2}{l}{24:00储电量} & 4348.26\\
\bottomrule\end{tabular}\end{table}

\FloatBarrier
\begin{table}[htbp]\centering\scriptsize
\caption{问题二 2025-06-21：指定时段及全天购电}
\begin{tabular}{lrrrrr}\toprule
时间段 & 购电量 & 时间段 & 购电量 & 时间段 & 购电量\\\midrule
10:00--10:10 & 0.00 & 12:00--12:10 & 0.00 & 14:00--14:10 & 0.00\\
16:00--16:10 & 250.70 & 18:00--18:10 & 0.00 & 20:00--20:10 & 0.00\\
\multicolumn{2}{l}{全天计划购电量} & 46473.41 & \multicolumn{2}{l}{全天实际购电费} & 26609.86\\
\bottomrule\end{tabular}\end{table}

\begin{table}[htbp]\centering\scriptsize
\caption{问题二 2025-06-21：四小时充放电与首末储电量}
\begin{tabular}{lrrrrr}\toprule
时间段 & 充电量 & 放电量 & 时间段 & 充电量 & 放电量\\\midrule
00:00--04:00 & 2869.46 & 0.00 & 04:00--08:00 & 833.33 & 2615.32\\
08:00--12:00 & 5087.51 & 0.00 & 12:00--16:00 & 3470.92 & 0.00\\
16:00--20:00 & 0.00 & 6017.99 & 20:00--24:00 & 8166.67 & 2622.01\\
\multicolumn{2}{l}{0:00储电量} & 2670.80 & \multicolumn{2}{l}{24:00储电量} & 8550.00\\
\bottomrule\end{tabular}\end{table}

\FloatBarrier
\begin{table}[htbp]\centering\scriptsize
\caption{问题二 2025-09-23：指定时段及全天购电}
\begin{tabular}{lrrrrr}\toprule
时间段 & 购电量 & 时间段 & 购电量 & 时间段 & 购电量\\\midrule
10:00--10:10 & 0.00 & 12:00--12:10 & 549.48 & 14:00--14:10 & 0.00\\
16:00--16:10 & 548.33 & 18:00--18:10 & 777.56 & 20:00--20:10 & 0.00\\
\multicolumn{2}{l}{全天计划购电量} & 70464.22 & \multicolumn{2}{l}{全天实际购电费} & 48722.41\\
\bottomrule\end{tabular}\end{table}

\begin{table}[htbp]\centering\scriptsize
\caption{问题二 2025-09-23：四小时充放电与首末储电量}
\begin{tabular}{lrrrrr}\toprule
时间段 & 充电量 & 放电量 & 时间段 & 充电量 & 放电量\\\midrule
00:00--04:00 & 1666.67 & 0.00 & 04:00--08:00 & 833.33 & 5644.61\\
08:00--12:00 & 6191.30 & 2995.39 & 12:00--16:00 & 4841.85 & 296.85\\
16:00--20:00 & 0.00 & 5695.38 & 20:00--24:00 & 8166.67 & 2944.62\\
\multicolumn{2}{l}{0:00储电量} & 8550.00 & \multicolumn{2}{l}{24:00储电量} & 8550.00\\
\bottomrule\end{tabular}\end{table}

\FloatBarrier
\begin{table}[htbp]\centering\scriptsize
\caption{问题二 2025-12-21：指定时段及全天购电}
\begin{tabular}{lrrrrr}\toprule
时间段 & 购电量 & 时间段 & 购电量 & 时间段 & 购电量\\\midrule
10:00--10:10 & 0.00 & 12:00--12:10 & 989.93 & 14:00--14:10 & 0.00\\
16:00--16:10 & 863.82 & 18:00--18:10 & 719.77 & 20:00--20:10 & 0.00\\
\multicolumn{2}{l}{全天计划购电量} & 98935.42 & \multicolumn{2}{l}{全天实际购电费} & 65505.96\\
\bottomrule\end{tabular}\end{table}

\begin{table}[htbp]\centering\scriptsize
\caption{问题二 2025-12-21：四小时充放电与首末储电量}
\begin{tabular}{lrrrrr}\toprule
时间段 & 充电量 & 放电量 & 时间段 & 充电量 & 放电量\\\midrule
00:00--04:00 & 1666.67 & 0.00 & 04:00--08:00 & 1666.67 & 1508.33\\
08:00--12:00 & 5833.33 & 7806.67 & 12:00--16:00 & 8241.15 & 2760.33\\
16:00--20:00 & 0.00 & 5609.57 & 20:00--24:00 & 8166.67 & 3030.43\\
\multicolumn{2}{l}{0:00储电量} & 8550.00 & \multicolumn{2}{l}{24:00储电量} & 8550.00\\
\bottomrule\end{tabular}\end{table}

\FloatBarrier
\begin{table}[htbp]\centering\scriptsize
\caption{问题二：指定日期全部紧急购电区间}
\setlength{\tabcolsep}{3pt}\begin{tabular}{lrlrlrlr}\toprule
\multicolumn{2}{c}{03-20} & \multicolumn{2}{c}{06-21} & \multicolumn{2}{c}{09-23} & \multicolumn{2}{c}{12-21}\\\midrule
时间段 & 电量 & 时间段 & 电量 & 时间段 & 电量 & 时间段 & 电量\\
00:20--00:30 & 12.39 & 00:20--00:30 & 40.44 & 00:30--01:20 & 22.17 & 01:00--01:20 & 29.05\\
01:10--01:40 & 37.48 & 01:20--01:30 & 21.10 & 02:00--02:20 & 55.41 & 02:10--02:20 & 0.61\\
03:00--03:30 & 70.17 & 01:50--02:10 & 13.24 & 03:50--04:00 & 11.95 & 03:00--03:10 & 15.38\\
04:50--05:10 & 48.10 & 02:50--03:30 & 99.52 & 04:50--05:10 & 34.28 & 03:30--03:50 & 108.42\\
05:50--06:00 & 7.85 & 18:50--19:10 & 78.17 & 05:40--05:50 & 10.72 & 05:50--06:00 & 8.04\\
08:50--09:10 & 26.14 & 20:20--20:30 & 15.05 & 06:20--06:30 & 1.78 & 06:10--06:30 & 41.11\\
11:10--11:20 & 8.27 & 22:00--22:30 & 71.19 & 10:00--10:20 & 13.05 & 10:30--11:30 & 271.92\\
12:30--12:50 & 26.69 & — & — & 10:50--11:00 & 5.98 & 15:20--15:30 & 13.09\\
15:30--16:20 & 208.66 & — & — & 14:30--15:20 & 272.38 & 18:40--19:00 & 8.15\\
16:30--17:20 & 61.84 & — & — & 15:50--16:50 & 147.40 & 19:40--20:10 & 63.76\\
18:40--19:10 & 73.67 & — & — & 17:30--17:40 & 16.92 & — & —\\
20:00--20:20 & 41.73 & — & — & 18:10--18:50 & 111.43 & — & —\\
21:10--21:30 & 42.68 & — & — & 19:10--19:30 & 45.49 & — & —\\
23:30--23:50 & 24.14 & — & — & 19:50--20:30 & 140.64 & — & —\\
— & — & — & — & 20:50--21:10 & 106.59 & — & —\\
— & — & — & — & 21:20--22:00 & 58.24 & — & —\\
— & — & — & — & 22:20--22:30 & 57.72 & — & —\\
— & — & — & — & 23:10--23:40 & 67.37 & — & —\\
合计 & 689.80 & 合计 & 338.71 & 合计 & 1179.52 & 合计 & 559.54\\
\bottomrule\end{tabular}\end{table}

\FloatBarrier
\subsection{问题三指定日期结果}
\begin{table}[htbp]\centering\scriptsize
\caption{问题三 2025-03-20：指定时段及全天购电}
\begin{tabular}{lrrrrr}\toprule
时间段 & 购电量 & 时间段 & 购电量 & 时间段 & 购电量\\\midrule
\multicolumn{6}{c}{0点原计划购电量}\\
10:00--10:10 & 0.00 & 12:00--12:10 & 551.93 & 14:00--14:10 & 0.00\\
16:00--16:10 & 544.34 & 18:00--18:10 & 705.08 & 20:00--20:10 & 0.00\\
\multicolumn{2}{l}{全天计划购电量} & 64668.95 & \multicolumn{2}{l}{全天实际购电费} & 43286.36\\
\multicolumn{6}{c}{最终调整后购电量（不含紧急购电）}\\
10:00--10:10 & 0.00 & 12:00--12:10 & 551.93 & 14:00--14:10 & 0.00\\
16:00--16:10 & 544.34 & 18:00--18:10 & 705.08 & 20:00--20:10 & 0.00\\
\multicolumn{5}{l}{全天最终调整后购电量} & 64101.47\\
\bottomrule\end{tabular}\end{table}

\begin{table}[htbp]\centering\scriptsize
\caption{问题三 2025-03-20：四小时充放电与首末储电量}
\begin{tabular}{lrrrrr}\toprule
时间段 & 充电量 & 放电量 & 时间段 & 充电量 & 放电量\\\midrule
00:00--04:00 & 1629.40 & 0.00 & 04:00--08:00 & 833.33 & 5973.45\\
08:00--12:00 & 6049.15 & 2666.55 & 12:00--16:00 & 5195.34 & 468.04\\
16:00--20:00 & 2.26 & 5675.55 & 20:00--24:00 & 3395.02 & 2974.39\\
\multicolumn{2}{l}{0:00储电量} & 8583.54 & \multicolumn{2}{l}{24:00储电量} & 4246.51\\
\bottomrule\end{tabular}\end{table}

\FloatBarrier
\begin{table}[htbp]\centering\scriptsize
\caption{问题三 2025-06-21：指定时段及全天购电}
\begin{tabular}{lrrrrr}\toprule
时间段 & 购电量 & 时间段 & 购电量 & 时间段 & 购电量\\\midrule
\multicolumn{6}{c}{0点原计划购电量}\\
10:00--10:10 & 0.00 & 12:00--12:10 & 0.00 & 14:00--14:10 & 0.00\\
16:00--16:10 & 223.91 & 18:00--18:10 & 0.00 & 20:00--20:10 & 0.00\\
\multicolumn{2}{l}{全天计划购电量} & 43570.65 & \multicolumn{2}{l}{全天实际购电费} & 25202.45\\
\multicolumn{6}{c}{最终调整后购电量（不含紧急购电）}\\
10:00--10:10 & 0.00 & 12:00--12:10 & 0.00 & 14:00--14:10 & 0.00\\
16:00--16:10 & 176.27 & 18:00--18:10 & 0.00 & 20:00--20:10 & 0.00\\
\multicolumn{5}{l}{全天最终调整后购电量} & 43321.19\\
\bottomrule\end{tabular}\end{table}

\begin{table}[htbp]\centering\scriptsize
\caption{问题三 2025-06-21：四小时充放电与首末储电量}
\begin{tabular}{lrrrrr}\toprule
时间段 & 充电量 & 放电量 & 时间段 & 充电量 & 放电量\\\midrule
00:00--04:00 & 762.83 & 0.00 & 04:00--08:00 & 857.88 & 2271.65\\
08:00--12:00 & 6553.33 & 0.00 & 12:00--16:00 & 4083.07 & 0.00\\
16:00--20:00 & 0.00 & 5995.04 & 20:00--24:00 & 8149.49 & 2645.72\\
\multicolumn{2}{l}{0:00储电量} & 2292.66 & \multicolumn{2}{l}{24:00储电量} & 8533.70\\
\bottomrule\end{tabular}\end{table}

\FloatBarrier
\begin{table}[htbp]\centering\scriptsize
\caption{问题三 2025-09-23：指定时段及全天购电}
\begin{tabular}{lrrrrr}\toprule
时间段 & 购电量 & 时间段 & 购电量 & 时间段 & 购电量\\\midrule
\multicolumn{6}{c}{0点原计划购电量}\\
10:00--10:10 & 0.00 & 12:00--12:10 & 502.56 & 14:00--14:10 & 0.00\\
16:00--16:10 & 560.29 & 18:00--18:10 & 767.73 & 20:00--20:10 & 0.00\\
\multicolumn{2}{l}{全天计划购电量} & 70219.95 & \multicolumn{2}{l}{全天实际购电费} & 47566.49\\
\multicolumn{6}{c}{最终调整后购电量（不含紧急购电）}\\
10:00--10:10 & 0.00 & 12:00--12:10 & 430.27 & 14:00--14:10 & 0.00\\
16:00--16:10 & 560.29 & 18:00--18:10 & 767.73 & 20:00--20:10 & 0.00\\
\multicolumn{5}{l}{全天最终调整后购电量} & 69537.44\\
\bottomrule\end{tabular}\end{table}

\begin{table}[htbp]\centering\scriptsize
\caption{问题三 2025-09-23：四小时充放电与首末储电量}
\begin{tabular}{lrrrrr}\toprule
时间段 & 充电量 & 放电量 & 时间段 & 充电量 & 放电量\\\midrule
00:00--04:00 & 1682.24 & 0.00 & 04:00--08:00 & 833.33 & 5844.03\\
08:00--12:00 & 6049.67 & 2598.82 & 12:00--16:00 & 4754.75 & 331.14\\
16:00--20:00 & 42.16 & 5706.45 & 20:00--24:00 & 8152.71 & 2962.87\\
\multicolumn{2}{l}{0:00储电量} & 8535.99 & \multicolumn{2}{l}{24:00储电量} & 8517.90\\
\bottomrule\end{tabular}\end{table}

\FloatBarrier
\begin{table}[htbp]\centering\scriptsize
\caption{问题三 2025-12-21：指定时段及全天购电}
\begin{tabular}{lrrrrr}\toprule
时间段 & 购电量 & 时间段 & 购电量 & 时间段 & 购电量\\\midrule
\multicolumn{6}{c}{0点原计划购电量}\\
10:00--10:10 & 0.00 & 12:00--12:10 & 965.84 & 14:00--14:10 & 0.00\\
16:00--16:10 & 872.53 & 18:00--18:10 & 708.71 & 20:00--20:10 & 0.00\\
\multicolumn{2}{l}{全天计划购电量} & 98380.06 & \multicolumn{2}{l}{全天实际购电费} & 65090.39\\
\multicolumn{6}{c}{最终调整后购电量（不含紧急购电）}\\
10:00--10:10 & 0.00 & 12:00--12:10 & 954.33 & 14:00--14:10 & 0.00\\
16:00--16:10 & 872.53 & 18:00--18:10 & 708.71 & 20:00--20:10 & 0.00\\
\multicolumn{5}{l}{全天最终调整后购电量} & 98056.59\\
\bottomrule\end{tabular}\end{table}

\begin{table}[htbp]\centering\scriptsize
\caption{问题三 2025-12-21：四小时充放电与首末储电量}
\begin{tabular}{lrrrrr}\toprule
时间段 & 充电量 & 放电量 & 时间段 & 充电量 & 放电量\\\midrule
00:00--04:00 & 1640.14 & 0.00 & 04:00--08:00 & 1660.82 & 1605.72\\
08:00--12:00 & 5884.07 & 7358.10 & 12:00--16:00 & 7766.93 & 2767.14\\
16:00--20:00 & 7.42 & 5602.70 & 20:00--24:00 & 8200.89 & 3043.38\\
\multicolumn{2}{l}{0:00储电量} & 8573.87 & \multicolumn{2}{l}{24:00储电量} & 8576.96\\
\bottomrule\end{tabular}\end{table}

\FloatBarrier
\begin{table}[htbp]\centering\scriptsize
\caption{问题三：指定日期全部紧急购电区间}
\setlength{\tabcolsep}{3pt}\begin{tabular}{lrlrlrlr}\toprule
\multicolumn{2}{c}{03-20} & \multicolumn{2}{c}{06-21} & \multicolumn{2}{c}{09-23} & \multicolumn{2}{c}{12-21}\\\midrule
时间段 & 电量 & 时间段 & 电量 & 时间段 & 电量 & 时间段 & 电量\\
00:20--00:30 & 12.39 & 00:20--00:30 & 40.44 & 00:30--01:20 & 22.17 & 01:00--01:20 & 29.05\\
01:10--01:40 & 37.48 & 01:20--01:30 & 21.10 & 02:00--02:20 & 55.41 & 02:10--02:20 & 0.61\\
03:00--03:30 & 70.17 & 01:50--02:10 & 13.24 & 03:50--04:00 & 11.95 & 03:00--03:10 & 15.38\\
04:50--05:10 & 48.10 & 02:50--03:30 & 99.52 & 04:50--05:10 & 34.31 & 03:30--03:50 & 108.42\\
05:50--06:00 & 6.72 & 18:50--19:10 & 79.29 & 05:40--05:50 & 17.70 & 05:50--06:00 & 8.04\\
08:10--09:20 & 147.65 & 20:20--20:30 & 18.51 & 06:10--06:30 & 0.57 & 06:20--06:30 & 15.63\\
09:40--10:00 & 56.01 & 22:00--22:30 & 70.63 & 07:40--07:50 & 5.03 & 10:30--11:40 & 363.36\\
11:00--11:30 & 114.30 & — & — & 14:20--15:00 & 75.47 & 13:00--13:10 & 1.78\\
11:40--12:00 & 35.27 & — & — & 15:50--16:40 & 126.89 & 18:40--19:00 & 6.49\\
12:30--12:50 & 109.16 & — & — & 17:30--17:50 & 19.82 & 19:40--20:00 & 59.05\\
13:10--13:20 & 0.81 & — & — & 18:10--18:50 & 96.87 & — & —\\
14:00--14:10 & 6.50 & — & — & 19:10--19:30 & 41.82 & — & —\\
15:50--16:10 & 22.60 & — & — & 19:50--20:30 & 131.14 & — & —\\
16:50--17:00 & 0.11 & — & — & 20:50--21:10 & 106.59 & — & —\\
18:40--19:10 & 85.04 & — & — & 21:20--22:00 & 78.37 & — & —\\
19:50--20:20 & 49.50 & — & — & 22:20--22:30 & 57.72 & — & —\\
21:10--21:30 & 47.58 & — & — & 23:10--23:40 & 67.37 & — & —\\
23:30--23:50 & 24.14 & — & — & — & — & — & —\\
合计 & 873.52 & 合计 & 342.74 & 合计 & 949.18 & 合计 & 607.81\\
\bottomrule\end{tabular}\end{table}

\FloatBarrier
\subsection{4-2指定日期结果}
\begin{table}[htbp]\centering\scriptsize
\caption{4-2 2025-03-20：指定时段及全天购电}
\begin{tabular}{lrrrrr}\toprule
时间段 & 购电量 & 时间段 & 购电量 & 时间段 & 购电量\\\midrule
10:00--10:10 & 0.00 & 12:00--12:10 & 653.15 & 14:00--14:10 & 12.32\\
16:00--16:10 & 494.35 & 18:00--18:10 & 0.00 & 20:00--20:10 & 731.74\\
\multicolumn{2}{l}{全天计划购电量} & 63984.26 & \multicolumn{2}{l}{全天实际购电费} & 43741.78\\
\bottomrule\end{tabular}\end{table}

\begin{table}[htbp]\centering\scriptsize
\caption{4-2 2025-03-20：四小时充放电与首末储电量}
\begin{tabular}{lrrrrr}\toprule
时间段 & 充电量 & 放电量 & 时间段 & 充电量 & 放电量\\\midrule
00:00--04:00 & 2500.00 & 0.00 & 04:00--08:00 & 1500.00 & 5904.30\\
08:00--12:00 & 5922.21 & 3144.53 & 12:00--16:00 & 5058.38 & 520.46\\
16:00--20:00 & 0.00 & 6405.44 & 20:00--24:00 & 833.33 & 2234.56\\
\multicolumn{2}{l}{0:00储电量} & 7950.00 & \multicolumn{2}{l}{24:00储电量} & 1950.00\\
\bottomrule\end{tabular}\end{table}

\FloatBarrier
\begin{table}[htbp]\centering\scriptsize
\caption{4-2 2025-06-21：指定时段及全天购电}
\begin{tabular}{lrrrrr}\toprule
时间段 & 购电量 & 时间段 & 购电量 & 时间段 & 购电量\\\midrule
10:00--10:10 & 0.00 & 12:00--12:10 & 0.00 & 14:00--14:10 & 0.00\\
16:00--16:10 & 252.87 & 18:00--18:10 & 487.95 & 20:00--20:10 & 0.00\\
\multicolumn{2}{l}{全天计划购电量} & 49515.72 & \multicolumn{2}{l}{全天实际购电费} & 24101.55\\
\bottomrule\end{tabular}\end{table}

\begin{table}[htbp]\centering\scriptsize
\caption{4-2 2025-06-21：四小时充放电与首末储电量}
\begin{tabular}{lrrrrr}\toprule
时间段 & 充电量 & 放电量 & 时间段 & 充电量 & 放电量\\\midrule
00:00--04:00 & 0.00 & 0.00 & 04:00--08:00 & 2749.81 & 2227.34\\
08:00--12:00 & 7535.00 & 0.00 & 12:00--16:00 & 3131.66 & 0.00\\
16:00--20:00 & 0.00 & 6017.99 & 20:00--24:00 & 9166.67 & 2757.01\\
\multicolumn{2}{l}{0:00储电量} & 1200.00 & \multicolumn{2}{l}{24:00储电量} & 9300.00\\
\bottomrule\end{tabular}\end{table}

\FloatBarrier
\begin{table}[htbp]\centering\scriptsize
\caption{4-2 2025-09-23：指定时段及全天购电}
\begin{tabular}{lrrrrr}\toprule
时间段 & 购电量 & 时间段 & 购电量 & 时间段 & 购电量\\\midrule
10:00--10:10 & 0.00 & 12:00--12:10 & 568.09 & 14:00--14:10 & 0.00\\
16:00--16:10 & 550.06 & 18:00--18:10 & 0.00 & 20:00--20:10 & 0.00\\
\multicolumn{2}{l}{全天计划购电量} & 70272.43 & \multicolumn{2}{l}{全天实际购电费} & 50538.08\\
\bottomrule\end{tabular}\end{table}

\begin{table}[htbp]\centering\scriptsize
\caption{4-2 2025-09-23：四小时充放电与首末储电量}
\begin{tabular}{lrrrrr}\toprule
时间段 & 充电量 & 放电量 & 时间段 & 充电量 & 放电量\\\midrule
00:00--04:00 & 4166.67 & 0.00 & 04:00--08:00 & 833.33 & 5614.15\\
08:00--12:00 & 6150.98 & 3025.85 & 12:00--16:00 & 4882.17 & 296.85\\
16:00--20:00 & 0.00 & 6245.22 & 20:00--24:00 & 4833.33 & 2394.78\\
\multicolumn{2}{l}{0:00储电量} & 6300.00 & \multicolumn{2}{l}{24:00储电量} & 5550.00\\
\bottomrule\end{tabular}\end{table}

\FloatBarrier
\begin{table}[htbp]\centering\scriptsize
\caption{4-2 2025-12-21：指定时段及全天购电}
\begin{tabular}{lrrrrr}\toprule
时间段 & 购电量 & 时间段 & 购电量 & 时间段 & 购电量\\\midrule
10:00--10:10 & 0.00 & 12:00--12:10 & 993.19 & 14:00--14:10 & 0.00\\
16:00--16:10 & 863.82 & 18:00--18:10 & 719.77 & 20:00--20:10 & 0.00\\
\multicolumn{2}{l}{全天计划购电量} & 97387.54 & \multicolumn{2}{l}{全天实际购电费} & 76456.54\\
\bottomrule\end{tabular}\end{table}

\begin{table}[htbp]\centering\scriptsize
\caption{4-2 2025-12-21：四小时充放电与首末储电量}
\begin{tabular}{lrrrrr}\toprule
时间段 & 充电量 & 放电量 & 时间段 & 充电量 & 放电量\\\midrule
00:00--04:00 & 666.67 & 0.00 & 04:00--08:00 & 0.00 & 833.33\\
08:00--12:00 & 5833.33 & 7806.67 & 12:00--16:00 & 8256.22 & 2772.54\\
16:00--20:00 & 0.00 & 5609.30 & 20:00--24:00 & 8333.33 & 3030.70\\
\multicolumn{2}{l}{0:00储电量} & 10200.00 & \multicolumn{2}{l}{24:00储电量} & 8700.00\\
\bottomrule\end{tabular}\end{table}

\FloatBarrier
\begin{table}[htbp]\centering\scriptsize
\caption{4-2：指定日期全部紧急购电区间}
\setlength{\tabcolsep}{3pt}\begin{tabular}{lrlrlrlr}\toprule
\multicolumn{2}{c}{03-20} & \multicolumn{2}{c}{06-21} & \multicolumn{2}{c}{09-23} & \multicolumn{2}{c}{12-21}\\\midrule
时间段 & 电量 & 时间段 & 电量 & 时间段 & 电量 & 时间段 & 电量\\
00:20--00:30 & 10.64 & 00:20--00:30 & 36.51 & 00:30--00:50 & 10.00 & 01:00--01:20 & 26.83\\
01:10--01:40 & 32.72 & 01:20--01:30 & 16.72 & 01:00--01:10 & 5.36 & 03:00--03:10 & 13.56\\
03:00--03:30 & 67.32 & 02:00--02:10 & 10.40 & 02:00--02:20 & 52.79 & 03:30--03:50 & 103.82\\
04:50--05:10 & 47.52 & 02:50--03:30 & 90.63 & 03:50--04:00 & 9.70 & 05:50--06:00 & 4.79\\
05:50--06:00 & 6.72 & 18:50--19:10 & 78.17 & 04:50--05:10 & 31.67 & 06:10--06:30 & 34.39\\
09:00--09:10 & 17.38 & 20:20--20:30 & 15.05 & 05:40--05:50 & 7.01 & 10:30--11:30 & 259.80\\
12:30--12:50 & 3.28 & 22:00--22:30 & 63.81 & 10:00--10:20 & 4.62 & 15:20--15:30 & 11.52\\
15:30--16:20 & 199.58 & — & — & 14:30--15:20 & 260.90 & 18:40--19:00 & 8.15\\
16:30--17:20 & 61.32 & — & — & 15:50--16:50 & 122.31 & 19:40--20:10 & 63.76\\
18:40--19:10 & 78.09 & — & — & 17:30--17:40 & 16.92 & — & —\\
20:00--20:20 & 43.60 & — & — & 18:10--18:50 & 96.79 & — & —\\
21:10--21:30 & 40.07 & — & — & 19:10--19:30 & 41.85 & — & —\\
23:30--23:50 & 23.20 & — & — & 19:50--20:30 & 140.17 & — & —\\
— & — & — & — & 20:50--21:10 & 105.01 & — & —\\
— & — & — & — & 21:20--22:00 & 54.90 & — & —\\
— & — & — & — & 22:20--22:30 & 57.72 & — & —\\
— & — & — & — & 23:10--23:40 & 46.50 & — & —\\
合计 & 631.45 & 合计 & 311.29 & 合计 & 1064.22 & 合计 & 526.64\\
\bottomrule\end{tabular}\end{table}

\FloatBarrier
\subsection{4-3指定日期结果}
\begin{table}[htbp]\centering\scriptsize
\caption{4-3 2025-03-20：指定时段及全天购电}
\begin{tabular}{lrrrrr}\toprule
时间段 & 购电量 & 时间段 & 购电量 & 时间段 & 购电量\\\midrule
\multicolumn{6}{c}{0点原计划购电量}\\
10:00--10:10 & 0.00 & 12:00--12:10 & 545.41 & 14:00--14:10 & 0.00\\
16:00--16:10 & 525.49 & 18:00--18:10 & -0.00 & 20:00--20:10 & 719.50\\
\multicolumn{2}{l}{全天计划购电量} & 67104.27 & \multicolumn{2}{l}{全天实际购电费} & 49343.99\\
\multicolumn{6}{c}{最终调整后购电量（不含紧急购电）}\\
10:00--10:10 & 0.00 & 12:00--12:10 & 545.41 & 14:00--14:10 & 0.00\\
16:00--16:10 & 525.49 & 18:00--18:10 & 0.00 & 20:00--20:10 & 719.50\\
\multicolumn{5}{l}{全天最终调整后购电量} & 66322.18\\
\bottomrule\end{tabular}\end{table}

\begin{table}[htbp]\centering\scriptsize
\caption{4-3 2025-03-20：四小时充放电与首末储电量}
\begin{tabular}{lrrrrr}\toprule
时间段 & 充电量 & 放电量 & 时间段 & 充电量 & 放电量\\\midrule
00:00--04:00 & 7500.00 & 0.00 & 04:00--08:00 & 3166.67 & 6274.00\\
08:00--12:00 & 5993.55 & 2366.00 & 12:00--16:00 & 5097.01 & 343.95\\
16:00--20:00 & 0.73 & 7170.36 & 20:00--24:00 & 65.10 & 1522.37\\
\multicolumn{2}{l}{0:00储电量} & 1200.00 & \multicolumn{2}{l}{24:00储电量} & 1200.00\\
\bottomrule\end{tabular}\end{table}

\FloatBarrier
\begin{table}[htbp]\centering\scriptsize
\caption{4-3 2025-06-21：指定时段及全天购电}
\begin{tabular}{lrrrrr}\toprule
时间段 & 购电量 & 时间段 & 购电量 & 时间段 & 购电量\\\midrule
\multicolumn{6}{c}{0点原计划购电量}\\
10:00--10:10 & 0.00 & 12:00--12:10 & 0.00 & 14:00--14:10 & 0.00\\
16:00--16:10 & 192.84 & 18:00--18:10 & 453.05 & 20:00--20:10 & 0.00\\
\multicolumn{2}{l}{全天计划购电量} & 34666.48 & \multicolumn{2}{l}{全天实际购电费} & 19165.30\\
\multicolumn{6}{c}{最终调整后购电量（不含紧急购电）}\\
10:00--10:10 & 0.00 & 12:00--12:10 & 0.00 & 14:00--14:10 & 0.00\\
16:00--16:10 & 134.07 & 18:00--18:10 & 453.05 & 20:00--20:10 & 0.00\\
\multicolumn{5}{l}{全天最终调整后购电量} & 34353.81\\
\bottomrule\end{tabular}\end{table}

\begin{table}[htbp]\centering\scriptsize
\caption{4-3 2025-06-21：四小时充放电与首末储电量}
\begin{tabular}{lrrrrr}\toprule
时间段 & 充电量 & 放电量 & 时间段 & 充电量 & 放电量\\\midrule
00:00--04:00 & 0.00 & 0.00 & 04:00--08:00 & 2586.28 & 2075.00\\
08:00--12:00 & 6553.33 & 0.00 & 12:00--16:00 & 4107.81 & 15.41\\
16:00--20:00 & 0.00 & 6096.96 & 20:00--24:00 & 53.86 & 2586.67\\
\multicolumn{2}{l}{0:00储电量} & 1200.00 & \multicolumn{2}{l}{24:00储电量} & 1200.00\\
\bottomrule\end{tabular}\end{table}

\FloatBarrier
\begin{table}[htbp]\centering\scriptsize
\caption{4-3 2025-09-23：指定时段及全天购电}
\begin{tabular}{lrrrrr}\toprule
时间段 & 购电量 & 时间段 & 购电量 & 时间段 & 购电量\\\midrule
\multicolumn{6}{c}{0点原计划购电量}\\
10:00--10:10 & 0.00 & 12:00--12:10 & 456.65 & 14:00--14:10 & 0.00\\
16:00--16:10 & 536.21 & 18:00--18:10 & 0.00 & 20:00--20:10 & 0.00\\
\multicolumn{2}{l}{全天计划购电量} & 66826.12 & \multicolumn{2}{l}{全天实际购电费} & 52706.02\\
\multicolumn{6}{c}{最终调整后购电量（不含紧急购电）}\\
10:00--10:10 & 0.00 & 12:00--12:10 & 306.14 & 14:00--14:10 & 0.00\\
16:00--16:10 & 536.21 & 18:00--18:10 & 0.00 & 20:00--20:10 & 0.00\\
\multicolumn{5}{l}{全天最终调整后购电量} & 66054.20\\
\bottomrule\end{tabular}\end{table}

\begin{table}[htbp]\centering\scriptsize
\caption{4-3 2025-09-23：四小时充放电与首末储电量}
\begin{tabular}{lrrrrr}\toprule
时间段 & 充电量 & 放电量 & 时间段 & 充电量 & 放电量\\\midrule
00:00--04:00 & 7333.33 & 0.00 & 04:00--08:00 & 3333.33 & 6441.96\\
08:00--12:00 & 6025.31 & 2198.04 & 12:00--16:00 & 4936.43 & 245.51\\
16:00--20:00 & 25.85 & 6307.45 & 20:00--24:00 & 16.57 & 2360.41\\
\multicolumn{2}{l}{0:00储电量} & 1200.00 & \multicolumn{2}{l}{24:00储电量} & 1200.00\\
\bottomrule\end{tabular}\end{table}

\FloatBarrier
\begin{table}[htbp]\centering\scriptsize
\caption{4-3 2025-12-21：指定时段及全天购电}
\begin{tabular}{lrrrrr}\toprule
时间段 & 购电量 & 时间段 & 购电量 & 时间段 & 购电量\\\midrule
\multicolumn{6}{c}{0点原计划购电量}\\
10:00--10:10 & 0.00 & 12:00--12:10 & 918.26 & 14:00--14:10 & 0.00\\
16:00--16:10 & 843.15 & 18:00--18:10 & 703.41 & 20:00--20:10 & 0.00\\
\multicolumn{2}{l}{全天计划购电量} & 95873.43 & \multicolumn{2}{l}{全天实际购电费} & 77661.40\\
\multicolumn{6}{c}{最终调整后购电量（不含紧急购电）}\\
10:00--10:10 & 0.00 & 12:00--12:10 & 918.26 & 14:00--14:10 & 0.00\\
16:00--16:10 & 844.01 & 18:00--18:10 & 703.41 & 20:00--20:10 & 0.00\\
\multicolumn{5}{l}{全天最终调整后购电量} & 95320.49\\
\bottomrule\end{tabular}\end{table}

\begin{table}[htbp]\centering\scriptsize
\caption{4-3 2025-12-21：四小时充放电与首末储电量}
\begin{tabular}{lrrrrr}\toprule
时间段 & 充电量 & 放电量 & 时间段 & 充电量 & 放电量\\\midrule
00:00--04:00 & 10666.67 & 0.00 & 04:00--08:00 & 3.60 & 864.86\\
08:00--12:00 & 5833.33 & 7351.09 & 12:00--16:00 & 7885.43 & 2899.17\\
16:00--20:00 & 12.56 & 5655.47 & 20:00--24:00 & 44.62 & 3030.85\\
\multicolumn{2}{l}{0:00储电量} & 1200.00 & \multicolumn{2}{l}{24:00储电量} & 1200.00\\
\bottomrule\end{tabular}\end{table}

\FloatBarrier
\begin{table}[htbp]\centering\scriptsize
\caption{4-3：指定日期全部紧急购电区间}
\setlength{\tabcolsep}{3pt}\begin{tabular}{lrlrlrlr}\toprule
\multicolumn{2}{c}{03-20} & \multicolumn{2}{c}{06-21} & \multicolumn{2}{c}{09-23} & \multicolumn{2}{c}{12-21}\\\midrule
时间段 & 电量 & 时间段 & 电量 & 时间段 & 电量 & 时间段 & 电量\\
00:20--00:30 & 22.46 & 00:10--00:30 & 74.72 & 00:30--01:20 & 111.02 & 00:00--00:20 & 7.92\\
01:00--02:00 & 82.55 & 01:20--02:20 & 123.00 & 01:40--02:30 & 112.37 & 01:00--01:20 & 57.62\\
02:30--02:40 & 0.58 & 02:50--03:30 & 149.65 & 03:40--04:00 & 34.84 & 01:30--01:40 & 11.18\\
03:00--03:40 & 111.77 & 03:50--04:10 & 13.90 & 04:50--05:20 & 67.28 & 02:10--02:20 & 15.66\\
04:50--05:20 & 81.18 & 17:20--17:40 & 14.78 & 05:30--05:50 & 29.41 & 02:40--02:50 & 0.88\\
05:50--06:00 & 20.46 & 18:50--19:10 & 105.08 & 06:10--06:30 & 28.01 & 03:00--03:20 & 37.48\\
06:40--07:00 & 7.10 & 20:20--20:40 & 52.87 & 06:50--07:10 & 12.98 & 03:30--03:50 & 144.79\\
08:00--09:20 & 342.10 & 21:40--22:40 & 125.64 & 07:40--08:00 & 26.75 & 04:00--04:20 & 11.40\\
09:40--10:00 & 80.43 & — & — & 09:30--09:50 & 23.79 & 05:40--06:30 & 97.47\\
10:20--10:30 & 17.89 & — & — & 10:00--10:20 & 7.35 & 07:10--07:20 & 3.74\\
10:50--12:00 & 461.66 & — & — & 10:40--11:00 & 35.51 & 07:50--08:00 & 2.03\\
12:30--13:30 & 191.60 & — & — & 11:30--11:50 & 22.71 & 08:30--09:00 & 24.22\\
14:00--14:20 & 35.44 & — & — & 12:50--13:00 & 0.16 & 09:30--09:50 & 39.29\\
14:50--15:00 & 7.47 & — & — & 13:50--14:00 & 14.31 & 10:20--11:50 & 574.75\\
15:30--16:20 & 90.12 & — & — & 14:20--15:20 & 217.15 & 12:00--12:20 & 39.04\\
16:40--17:00 & 27.47 & — & — & 15:50--16:50 & 263.19 & 12:40--13:10 & 63.72\\
18:40--19:20 & 127.95 & — & — & 17:00--17:50 & 73.03 & 15:20--15:30 & 5.86\\
19:30--20:20 & 93.09 & — & — & 18:00--18:50 & 195.57 & 16:20--16:30 & 11.34\\
21:10--21:40 & 75.12 & — & — & 19:10--19:30 & 82.27 & 18:00--18:20 & 11.27\\
23:30--23:50 & 60.55 & — & — & 19:50--22:00 & 562.42 & 18:40--19:00 & 30.66\\
— & — & — & — & 22:10--22:40 & 77.99 & 19:40--20:10 & 88.22\\
— & — & — & — & 23:10--23:40 & 123.17 & 21:50--22:00 & 2.84\\
— & — & — & — & — & — & 22:50--23:00 & 12.06\\
— & — & — & — & — & — & 23:30--23:40 & 2.50\\
合计 & 1936.99 & 合计 & 659.63 & 合计 & 2121.28 & 合计 & 1295.96\\
\bottomrule\end{tabular}\end{table}

\FloatBarrier

\FloatBarrier
\section{支撑材料与复现说明}
本文主结果对应\texttt{result1.xlsx}、\texttt{result2.xlsx}、\texttt{result3.xlsx}、\texttt{result4-2.xlsx}和\texttt{result4-3.xlsx}，分别位于各问正式结果目录。它们包含题面要求的完整策略，本文没有覆盖或重算这些文件。

论文图片和指定结果表由独立资料准备脚本读取已有结果生成；数据映射与SHA-256记录见本稿目录的\texttt{source\_manifest.json}。该脚本仅读取题目、指定主模型相关文件、新图包和效率补充材料，不调用优化器。LaTeX编译与资产校验过程在本稿说明文件中记录。

主模型的计算可依据各问运行说明复现。第二问正式方案需使用有效残差起点7和56天等权；第三问及4-3的四次更新是预设主方案。效率对照材料单独列示来源，不将其视为本次重新完成的实验。参考资料仅列本文实际使用的题面及内部研究材料，尚未加入未经查阅的外部方法文献。

\end{document}
